% มคอ.5 — รายงานการสัมฤทธิ์ผล 7015101 ภาคเรียนที่ 2/2568
% compile: xelatex AUN3_7015101.tex
\documentclass{mko5}

\begin{document}

\mkofivetitle

\coursecode{7015101}
\coursename{เทคโนโลยีอุบัติใหม่ทางวิศวกรรมคอมพิวเตอร์และปัญญาประดิษฐ์}
\curriculumnote{หลักสูตรวิศวกรรมศาสตรมหาบัณฑิต สาขาวิชาวิศวกรรมคอมพิวเตอร์และปัญญาประดิษฐ์ พ.ศ. 2568 (แผนสารนิพนธ์)}
\coursegroup{วิชาเฉพาะด้านบังคับ}
\coursecredits{3 (2-2-5)}
\studyplan{แผน 2 --- วิชาชีพ (สารนิพนธ์)}
\semester{ภาคเรียนที่ 2 / ปีการศึกษา 2568}
\prerequisite{ไม่มี}
\corequisite{ไม่มี}
\studyplace{มหาวิทยาลัยราชภัฏอุตรดิตถ์ --- คณะเทคโนโลยีอุตสาหกรรม / ห้องเรียนและห้องปฏิบัติการ ICIT}
\registeredstudents{4 คน}
\completedstudents{4 คน}
\instructor{ผศ.ดร.พิศิษฐ์ นาคใจ และ ผศ.ดร.พัชรี มณีรัตน์}
\coinstructors{---}

\mkoheaderinfo

% ----- 1. กิจกรรมที่ดำเนินการจริง (จาก มคอ.3 + บันทึกประชุม 27 มี.ค. 2569) -----
\begin{teachingactualtable}
  \actualrow{CLO1}{บรรยายแบบมีปฏิสัมพันธ์ นำเสนอกรณีศึกษา (AI/ML, Blockchain, Quantum) อภิปรายกลุ่ม
  ครอบคลุม IoT, Edge, Cloud-native ตามแผนรายสัปดาห์ สัปดาห์ 1--9}
  \actualrow{CLO2}{สาธิตและ Workshop (IoT, Cloud, DevOps) Project-based Learning
  สัปดาห์ที่ 10: IEL Community Workshop (26 มี.ค. 2569) นักศึกษาถ่ายทอดความรู้ AI ให้โรงเรียน 3 แห่งใน จ.อุตรดิตถ์
  พัฒนาโครงงานสัปดาห์ 11--14 นำเสนอผลงานสัปดาห์ 15}
  \actualrow{CLO3}{เชิญวิทยากร อภิปรายผลกระทบต่อสังคมและจริยธรรมเทคโนโลยี
  บทสะท้อนคิดและรายงานกลางภาค}
  \actualrow{CLO4}{กิจกรรมแลกเปลี่ยนความคิดเห็น การมีส่วนร่วมในชั้นเรียนและ Workshop ชุมชน
  การนำเสนอโครงงานและสรุปรายวิชา}
  \actualrow{CLO5}{สาธิตขั้นตอนการใช้งาน Lab/Cloud/GPU ให้คำแนะนำระหว่างปฏิบัติ
  ทดสอบและสาธิตระบบโครงงาน การประเมินปฏิบัติและสังเกตพฤติกรรม}
\end{teachingactualtable}

% ----- 2. วิธีการวัดผล (จาก มคอ.3 ข้อ 4.3) -----
\begin{assessmentusedtable}
  \assessusedrow{CLO1}{การสอบข้อเขียน (สอบกลางภาค) การนำเสนอ รายงานการวิเคราะห์เทคโนโลยี --- Rubric 4 ระดับ}
  \assessusedrow{CLO2}{สอบกลางภาค ผลงานโครงงาน การสอบปฏิบัติ การนำเสนอผลงานและรายงาน IEL Workshop}
  \assessusedrow{CLO3}{การมีส่วนร่วมในชั้นเรียน บทความสะท้อนคิด รายงานกลางภาค}
  \assessusedrow{CLO4}{การมีส่วนร่วมในชั้นเรียน บทสะท้อนคิด การนำเสนอโครงงาน}
  \assessusedrow{CLO5}{การทดสอบปฏิบัติ การสังเกตพฤติกรรม สอบปลายภาค การประเมินโครงงาน}
\end{assessmentusedtable}

% ----- 3. สรุปสัมฤทธิ์ผล (n=4, ผ่านเกณฑ์ระดับ 3 ขึ้นไปทุก CLO --- บันทึกประชุม 5/ภาค 2/2568) -----
\mkosectionthree

\closummaryblock{CLO1}{อธิบาย แนวคิดและหลักการของเทคโนโลยีอุบัติใหม่ทางวิศวกรรมคอมพิวเตอร์และปัญญาประดิษฐ์ ได้อย่างถูกต้อง}{80}
\achievementpass{4}{100}{ผ่านเกณฑ์ Rubric ระดับ 3 ขึ้นไปทุกคน (สอบกลางภาคและรายงานผ่าน)}
\achievementfail{0}{0}{ไม่มี}
\closummaryend

\closummaryblock{CLO2}{ประยุกต์ใช้ เทคโนโลยีอุบัติใหม่ ในการแก้ปัญหาทางวิศวกรรมคอมพิวเตอร์และปัญญาประดิษฐ์}{80}
\achievementpass{4}{100}{ผ่านโครงงาน Workshop ชุมชน และสอบปลายภาค}
\achievementfail{0}{0}{ไม่มี}
\closummaryend

\closummaryblock{CLO3}{รับรู้ ผลกระทบและการเปลี่ยนแปลงของเทคโนโลยีอุบัติใหม่ ที่มีต่อสังคมและวิชาชีพ}{80}
\achievementpass{4}{100}{ผ่านการมีส่วนร่วมและบทสะท้อนคิด}
\achievementfail{0}{0}{ไม่มี}
\closummaryend

\closummaryblock{CLO4}{ตอบสนอง ต่อการเรียนรู้เทคโนโลยีใหม่ ด้วยความกระตือรือร้นและมีส่วนร่วม}{80}
\achievementpass{4}{100}{เข้าร่วม IEL Workshop และนำเสนอโครงงานครบ}
\achievementfail{0}{0}{ไม่มี}
\closummaryend

\closummaryblock{CLO5}{ปฏิบัติตาม แนวทางการใช้เทคโนโลยีอุบัติใหม่ ภายใต้การแนะนำได้อย่างเหมาะสม}{80}
\achievementpass{4}{100}{ผ่านการทดสอบปฏิบัติและประเมินโครงงาน}
\achievementfail{0}{0}{ไม่มี}
\closummaryend

% ----- 4--6 -----
\mkosectionfour
\begin{mkotextblock}
จำนวนนักศึกษาน้อย ($n=4$) ทำให้การวิเคราะห์เชิงสถิติมีข้อจำกัด;
การจัด IEL Community Workshop ต้องประสานกับโรงเรียนและเดินทางนอกสถานที่ในช่วงสัปดาห์ที่ 10
\end{mkotextblock}

\mkosectionfive
\begin{mkotextblock}
เพิ่มชุดข้อมูลฝึกปฏิบัติและตัวอย่างโครงงานจากภาคอุตสาหกรรม;
จัดทำแบบฟอร์มสะท้อนคิดมาตรฐานสำหรับ CLO3--CLO4;
เชื่อมโยงผล มคอ.5 กับการทบทวนระดับหลักสูตรในภาคเรียนที่ 3/2568
\end{mkotextblock}

\mkosectionsix
\begin{mkotextblock}
ขอให้คณะสนับสนุนการเข้าถึง GPU/Cloud สำหรับโครงงาน CPE\&AI;
พิจารณาเพิ่มคำถามชี้นำในการประเมิน Oral Presentation เพื่อตรวจสอบความเข้าใจรายบุคคล
(ตามข้อเสนอจากการประชุมคณะกรรมการบริหารหลักสูตร ครั้งที่ 1/2569)
\end{mkotextblock}

\mkosectionseven
\mkosignaturepair{ผศ.ดร.พิศิษฐ์ นาคใจ และ ผศ.ดร.พัชรี มณีรัตน์}{30 เมษายน 2569}{ผศ.ดร.พิทักษ์ คล้ายชม}{30 เมษายน 2569}

\end{document}
